% ============================================================
% PAPER 5 of 11 -- exact_belief_computation
% Assembled 2026-09-29 by content-and-merits partition.
% Title: Exact Belief Computation on the Observation Ladder
% Source: arena agent 1/paper rewrites/latex/paper2_exact_belief_computation_v13.tex
% Partition note: Single-source. Short self-contained result: Hamming classification and the antichain census.
%% STATUS: structurally assembled as a NEW version. Editorial integration
%%         (de-duplication of shared framework, sibling cross-citations)
%%         is NOT yet done. See PAPERS_MANIFEST.md.
%% ============================================================

%% v12 (2026-09-29): (a) SCOPE DECISION RESOLVED -- expand, do not fold into paper 3.
%%   Paper 3 CITES this paper as "the exact-computation companion"; folding would join
%%   two different mathematical cores and break that citation.
%% (b) 'ebc' identification CONFIRMED, moderate->high: 16 Lean modules named EBC_*
%%   (Bands/Classification/Deadline/Dynamics/ExactBelief/Hamming/Ladder/Pairs) in
%%   namespace Formalizations.EBC map onto this paper's own sections.
%% (c) 'Scale II' relation RESOLVED: the parent theory (belief-state safety value
%%   V_k(b)) is paper2_belief_state_v2 -- paper 2 of this set in its later form,
%%   32.6%% of its sentences appear verbatim in paper02. Stated in the text.
%% (d) Added subsection priorart: concedes Sperner/Dedekind/Kleitman and Hamming
%%   adjacency as classical, states this paper proves NO new theorem about
%%   antichains, narrows the claim to four items. Added subsection mech (16 EBC
%%   modules, 134 declarations, 0 sorry) with the honest note that the census
%%   enumeration itself is NOT formalized.
%% (e) Markups are \subsection to match this paper (it uses no \section).
%% PRIOR ART AND SCOPE REQUIRED -- bar item, inserted 2026-09-29.
%% This is the shortest paper in the family (~3.6k words) with no related-work
%% position, and it is the furthest from "substantial methodological development".
%% Two things must happen before it can stand:
%%   (i) PRIOR ART: exact and rational-arithmetic computation of viability
%%       objects -- Saint-Pierre's kernel algorithms and their successors -- and
%%       the combinatorial content, i.e. Sperner and antichain bounds on the
%%       observation ladder. The antichain census needs a comparator.
%%   (ii) SCOPE: at this length it is a short-format result. Either expand it
%%       beyond the cube instance to the general case, or fold it into Paper 3,
%%       which is the other computational paper. Decide before submission; do not
%%       submit at this length with neither done.

% SIBLING PAPERS: Paper 1 (Obstruction Certificates under Incomplete Observation), Paper 2 (Exact Probabilistic Sufficiency), Paper 3 (Computational Certification of Obstruction), Paper 4 (Minimax Dual Certificates).
% Cross-cite, do not re-derive: the selector principle, the epistemic kernel
% and the observation structure are defined in Paper 1. Papers 2-5 cite it.

% SIBLING PAPERS: Paper 1 (Obstruction Certificates under Incomplete Observation), Paper 2 (Exact Probabilistic Sufficiency), Paper 3 (Computational Certification of Obstruction), Paper 4 (Minimax Dual Certificates).
% Cross-cite, do not re-derive: the selector principle, the epistemic kernel
% and the observation structure are defined in Paper 1. Papers 2-5 cite it.

% RESOURCED 2026-09-29: previous assembly (v10) was built from _v9 of this source,
%   4 versions behind. Re-assembled against paper2_exact_belief_computation_v13.tex. Content is a superset:
%   no labels dropped. See VERSION_CURRENCY.md.

\documentclass[10pt,twocolumn]{article}
\usepackage[a4paper,margin=15mm]{geometry}
\usepackage{amsmath,amssymb}
\usepackage{amsthm}
\theoremstyle{plain}
\newtheorem{theorem}{Theorem}
\newtheorem{proposition}[theorem]{Proposition}
\newtheorem{lemma}[theorem]{Lemma}
\newtheorem{corollary}[theorem]{Corollary}
\theoremstyle{definition}
\newtheorem{definition}[theorem]{Definition}
\theoremstyle{remark}
\newtheorem{remark}[theorem]{Remark}
\usepackage{booktabs,array,calc}
\usepackage[font=small,labelfont=bf]{caption}
\usepackage[colorlinks=true]{hyperref}
\providecommand{\tightlist}{\setlength{\itemsep}{0pt}\setlength{\parskip}{0pt}}
\emergencystretch=3em
\allowdisplaybreaks

\begin{document}
\title{Exact Belief-State Computation at Scale II:\\
The Four-Parameter Cube and the Antichain Census\\[0.2em]
{\large\rmfamily A computational companion to the belief-state
safety-value theory}}
\author{Amin Abaee\\[0.35em]
{\small Independent Researcher}\\[0.55em]
{\small\href{https://orcid.org/0000-0002-0019-1842}{ORCID: 0000-0002-0019-1842}}\\[0.3em]
{\small\href{mailto:amin\_abaee@ut.ac.ir}{amin\_abaee@ut.ac.ir}}}
\date{September 26, 2026}
\maketitle

\begin{abstract}
This paper pushes the exact discipline of belief-state safety values into the curse-of-dimensionality regime: a four-parameter cube with
\(2^{4} = 16\) hidden-parameter cells on \(16\) stock levels ---
\(256\) augmented cells, \(17\) controls, every computation in exact
integer and rational arithmetic. The structure is classified
completely. A pairwise-Hamming lemma (every pair of distinct actions
sums, over any two parameter cells, to at most twice their agreement
count) forces every blind survivable set to consist of
Hamming-adjacent cells, and no three cells of the four-cube are
pairwise adjacent: no blind policy of any period keeps three or more
cells safe. The bands are exact: \(16\) maximal singletons at the
floor's edge, exactly the \(32\) Hamming-adjacent pairs at every level
above it (the matched--singly-mismatched alternation rises at net
\(\tfrac15\) per two steps with a \(\tfrac1{10}\) dip). The
observation ladder doubles: each additional exact probe of one
parameter doubles the conditional kernel mass ---
\(\tfrac1{16} \to \tfrac18 \to \tfrac14 \to \tfrac12 \to 1\) at the
edge, \(\tfrac18 \to \tfrac14 \to \tfrac12 \to 1 \to 1\) above it ---
so three probes exhaust observation above the edge. Exact point-based
evaluation is verified against definition-level enumeration at seven
rational beliefs, three levels, four horizons; the four-step alpha-set
deduplicates \(83{,}521\) action sequences to at most \(545\) distinct
outcome vectors. The antichain census is the headline:
\(1{,}048{,}576\) raw subset evaluations collapse to \(496\) stored
maximal sets. A third instance of the theory's additive deadline law closes the paper (threshold \(1 + T/2\), zero probe
excursion).
\end{abstract}

\textbf{Keywords:} exact computation; partially observable control;
zero-suppressed sets; survivable sets; Hamming structure; observation
design; antichain compression

\subsection{Introduction}\label{introduction}

Two working disciplines govern the exact treatment of belief-state safety values: evaluate exactly where you probe, and store the antichain of maximal survivable sets rather than the lattice (Abaee, 2026, An obstruction calculus). This paper asks whether the exact discipline survives the curse of dimensionality, and what new structure appears there. The answer is affirmative, and the structure changes: the four-parameter cube is not a larger version of the two-parameter instance but a different regime. The cube is the paper's running instance; every structural claim below is instantiated on it.
Its survivable sets are governed by Hamming adjacency rather than by
coordinate matching; its observation ladder is geometric rather than
one-shot; and its raw subset lattice (\(2^{16}\) per level) is four
orders of magnitude larger than its stored antichain. All claims below
are regenerated by the paper's verification script in exact integer
and rational arithmetic, standard library only. The contributions are:
the pairwise-Hamming classification of the cube's survivable-set
structure (Section \ref{hamming}); the exact observation-ladder
doubling and the point-based evaluation at four parameters (Sections
\ref{ladder}--\ref{census}); and the antichain census with the
audited-structure classification and the deadline instance (Sections
\ref{glance}--\ref{deadline}).

\subsection{The cube instance}\label{instance}

The cube's local vocabulary (\(z\), \(z_{0}\), \(T\), the cell
structure) is disjoint from the companion papers' colliding letters;
the horizon \(T\) carries the same meaning as the certification
companion's \(J_{\mathcal{I}}(T)\). The state is \((z, \theta)\) with stock \(z\) on the grid
\(z_{0} \in \{1.0, 1.1, \dots, 2.5\}\) and hidden parameter
\(\theta = (\theta_{1}, \theta_{2}, \theta_{3}, \theta_{4}) \in
\{+1, -1\}^{4}\): \(16\) parameter cells on \(16\) levels, \(256\)
augmented cells. Controls are \(u \in \{+1,-1\}^{4} \cup \{(0,0,0,0)\}\)
(\(17\) actions), and the one-step drift is
\[
\dot z \;=\; d(u, \theta) \;=\; -\tfrac12
\;+\; \tfrac15\,\textstyle\sum_{i=1}^{4} u_{i}\theta_{i}:
\]
the matched action \(u = \theta\) rises at \(+\tfrac3{10}\); a single
mismatch gives \(-\tfrac1{10}\), two give \(-\tfrac12\), three give
\(-\tfrac9{10}\), the antimatched action \(-\tfrac{13}{10}\), and the
hold action \(-\tfrac12\) for every cell. The floor is \(z \ge 1\).
The sum over all sixteen cells of any action's drifts is \(-8\) per
step: the cube as a whole always declines, and the blind controller
declared here sees no observation of \(\theta\).

\subsection{The pairwise-Hamming classification}\label{hamming}

\begin{lemma}[pair-sum bound and its survival cost]\label{lem:pairsum}
Fix the dimension \(m \ge 1\), cells \(\theta, \theta' \in
\{+1,-1\}^{m}\) at Hamming distance \(h\), the action set
\(U = \{+1,-1\}^{m} \cup \{(0,\dots,0)\}\), the drift
\(d(u, \theta) = -\tfrac12 + \tfrac15\langle u, \theta\rangle\),
the floor \(z \ge 1\), and the blind class: one action sequence for
every parameter branch, of any period and possibly aperiodic. Then:
(i) for every \(u \in \{+1,-1\}^{m}\),
\(\langle u, \theta\rangle + \langle u, \theta'\rangle
\le 2\,(m - h)\), and the hold action gives \(0\) for every pair;
(ii) if a blind policy keeps both cells above the floor indefinitely,
then the pair's long-run average inner product satisfies
\(\liminf_{T \to \infty} \frac1T \sum_{t < T}
\bigl(\langle u_{t}, \theta\rangle + \langle u_{t},
\theta'\rangle\bigr) \ge 5\).
\end{lemma}

\begin{proof}
(i) The coordinates split into \(m - h\) agreements and \(h\)
disagreements; on an agreement \(u_{i}\theta_{i} = u_{i}\theta'_{i}
\le 1\) and on a disagreement the two terms cancel, so the sum is
\(2 \sum_{\text{agree}} u_{i}\theta_{i} \le 2(m-h)\); the hold
action gives \(0\) termwise. (ii) Survival of cell \(\theta\) means
\(z_{0} + \sum_{t < T} d(u_{t}, \theta) \ge 1\) for every \(T\);
dividing by \(T\) and passing to the limit, the Ces\`aro average of
the drifts has \(\liminf \ge 0\), and since
\(d(u_{t}, \theta) = -\tfrac12 + \tfrac15\langle u_{t},
\theta\rangle\), the \(\liminf\) of the average of
\(\langle u_{t}, \theta\rangle\) is at least
\(\tfrac52\); the same holds for \(\theta'\), and the two
\(\liminf\)s add. For a periodic policy the average is the exact
cycle average and no \(\liminf\) is needed.
\end{proof}

\begin{corollary}[the cube's pairs are Hamming-adjacent]
\label{cor:hamming}
At \(m = 4\): (ii) forces the pair average to reach \(5\) while (i)
caps it at \(2(4 - h)\), so \(5 \le 2(4 - h)\), i.e.\
\(h \le \tfrac32\): every blind survivable pair has \(h \le 1\),
and since survivability passes to subsets, every pair inside a blind
survivable set is Hamming-adjacent. In general the argument gives
\(h \le m - \tfrac52\) (Remark~\ref{rem:scope}).
\end{corollary}

\begin{lemma}[no adjacency triangles]\label{lem:triangle}
No three distinct cells of \(\{+1,-1\}^{m}\) (any \(m\)) are
pairwise at Hamming distance \(1\). Hence every blind survivable set
at \(m = 4\) has at most two members, under every blind policy ---
any period, and no period at all. The bound (i) of
Lemma~\ref{lem:pairsum} is verified for all \(120\) cell pairs and
all \(17\) actions at \(m = 4\) (and for all \(496\) pairs and
\(33\) actions at \(m = 5\)), and
Lemma~\ref{lem:triangle} over all \(\binom{16}{3} = 560\) triples, by the verification script.
\end{lemma}

\begin{proof}
Flip coordinates so that the first cell is \(\mathbf{1}\) (flips
preserve all Hamming distances). A second cell at distance \(1\) is
\(\mathbf{1}\) with exactly one coordinate \(i\) flipped. A third
distinct cell at distance \(1\) from \(\mathbf{1}\) is
\(\mathbf{1}\) with exactly one coordinate \(j\) flipped: if
\(j = i\) it is the second cell, and if \(j \neq i\) the two differ
from each other in the two coordinates \(i, j\), so their Hamming
distance is \(2\). No triangle exists. (Equivalently: the hypercube
is bipartite by coordinate parity, so it has no odd cycle.)
\end{proof}

\begin{proposition}[bands and maximality]\label{prop:bands}
A singleton \(\{\theta\}\) is survivable from every \(z_{0} \ge 1\)
(matched play rises at \(+\tfrac3{10}\)). A Hamming-adjacent pair
\(\{\theta, \theta'\}\) is survivable exactly from
\(z_{0} \ge 1.1\): the alternation \((\theta, \theta')\) gives each
member drifts \(+\tfrac3{10}, -\tfrac1{10}\) per two steps --- net
\(+\tfrac15\) with a within-cycle dip of \(\tfrac1{10}\) --- certified
at \(60\) steps, and from \(z_{0} = 1.0\) the first cell to act under
its mismatched action dips below the floor. The maximal survivable
sets are therefore exactly the \(16\) singletons at \(z_{0} = 1.0\)
and exactly the \(32\) Hamming-adjacent pairs at every
\(z_{0} \ge 1.1\), and nothing larger under any blind policy. An
exhaustive search over all \(5{,}219\) one-, two-, and three-periodic
policies with \(60\)-step certificates reproduces this classification
independently.
\end{proposition}

\begin{proof}
Singletons: matched play has drift \(+\tfrac3{10}\) from every
\(z_{0} \ge 1\). Pairs: let \(\{\theta, \theta'\}\) be
Hamming-adjacent and alternate \((\theta, \theta')\); each member
takes the matched drift \(+\tfrac3{10}\) and the mismatched drift
\(-\tfrac1{10}\) once per cycle, a net of \(+\tfrac15\) per two
steps with the cycle minimum \(z_{0} - \tfrac1{10}\) after the
mismatched step --- survivable exactly from \(z_{0} \ge 1.1\)
(the \(60\)-step certificates confirm the cycle arithmetic).
No pair from \(z_{0} = 1.0\): for any first action \(u\), if
\(u = \theta\) then \(\langle u, \theta'\rangle = 2\) and cell
\(\theta'\) falls to \(z_{0} - \tfrac1{10} < 1\) in one step, while
any other \(u\) drives at least one member below the floor in its
first step (the hold gives \(-\tfrac12\) to both); no repair exists
because the floor binds at every step. Nothing larger: any set of
three or more cells contains two at Hamming distance \(\ge 2\)
(Lemma~\ref{lem:triangle}), and that pair violates
Corollary~\ref{cor:hamming}, so no blind policy keeps it above the
floor --- and survivability passes to subsets, so the whole set fails
with it. Maximality is immediate: singletons extend to pairs from
\(1.1\), pairs extend to nothing.
\end{proof}

\begin{remark}[scope in dimension: the ball construction]
\label{rem:scope}
The pairs-only census is the four-dimensional cube's, and the boundary
is sharp. In general dimension the pair argument of
Corollary~\ref{cor:hamming} gives only \(h \le m - \tfrac52\); and
the constant matched action \(u = \mathbf{1}\) gives a cell at
Hamming distance \(k\) the constant drift
\(-\tfrac12 + \tfrac{m - 2k}{5}\), nonnegative exactly when
\(k \le \tfrac{2m-5}{4}\). Hence the Hamming ball of radius
\(r^{*}(m) = \lfloor (2m - 5)/4 \rfloor\) around any cell is
blind-survivable from every \(z_{0} \ge 1\):
\(r^{*}(4) = 0\), \(r^{*}(5) = r^{*}(6) = 1\),
\(r^{*}(7) = r^{*}(8) = 2\). At \(m = 5\) the radius-one ball ---
\(1 + 5 = 6\) cells, drifts \(+\tfrac12\) and \(+\tfrac1{10}\)
under constant matched play --- survives from the floor's edge, where
the four-cube keeps only singletons: the classification changes
qualitatively at \(m = 5\), and the maximal sets for \(m \ge 5\)
are not classified here. All eight radii and the sign pattern are
verified by the verification script, with \(60\)-step simulations of the
radius-one ball at \(m = 5\) and the \(29\)-cell radius-two ball at
\(m = 7\).
\end{remark}

\begin{remark}[the crude instrument is not sufficient]
\label{rem:crude}
The averaging identity used on smaller instances --- a survivable set
\(S\) needs \(\sum_{i} |\sum_{\theta \in S} \theta_{i}| \ge
\tfrac52\,|S|\) --- is not sufficient here: the four-set
\(\{++++, +++-, ++-+, +-++\}\) attains the bound with equality
(\(4 + 2 + 2 + 2 = 10\)) yet is not survivable, since it pairs cells
at Hamming distance \(2\). The pairwise instrument of
Lemma~\ref{lem:pairsum} is the sharp one.
\end{remark}

\subsection{The observation ladder doubles}\label{ladder}

\begin{proposition}[geometric ladder]\label{prop:ladder}
Prior \(b_{0}\) is uniform on the sixteen cells; a probe of one
parameter refines the support to a subcube. The maximal conditional
kernel mass doubles with each probe:
\begin{align*}
\tfrac1{16} \to \tfrac18 \to \tfrac14 \to \tfrac12 \to 1
&\quad (z_{0} = 1.0),\\
\tfrac18 \to \tfrac14 \to \tfrac12 \to 1 \to 1
&\quad (z_{0} \ge 1.1),
\end{align*}
for zero, one, two, three, and four probes. Above the band's edge
three probes exhaust observation --- the support is then a
Hamming-adjacent pair, jointly survivable, and the fourth probe buys
nothing --- while at the edge the fourth probe is worth
\(\tfrac12\), the single-cell mass the pair's dip excludes.
\end{proposition}

\begin{proof}
The companion theory's antichain formula makes the kernel value of a
prior the maximal survivable-set mass it charges, and survivability is
hereditary downward --- a policy keeping a set alive keeps every
subset alive --- so the value at the uniform prior on a support
\(S\) is \(\max_{A \subseteq S,\ A
\text{ survivable}} |A|/|S|\), attained inside the support. By
Proposition~\ref{prop:bands} the survivable subsets of a subcube are
singletons and Hamming-adjacent pairs, so a \(2^{j}\)-cell subcube
has value \(2/2^{j}\) from \(z_{0} \ge 1.1\) and \(1/2^{j}\) at
the edge. The supports after \(0, 1, 2, 3, 4\) probes of one
parameter are subcubes with \(j = 4, 3, 2, 1, 0\), giving the two
ladders exactly; the fourth edge probe buys the last cell because at
the edge only singletons survive and the fully probed support is then
itself the surviving cell. The script recomputes all ten values by
exhaustive survivable-subset enumeration within each support.
\end{proof}

\subsection{Exact point-based evaluation at four parameters}
\label{pbvi}

The evaluation is point-based in the POMDP literature's sense
(Lovejoy, 1991), with the distinction that the values here are exact at
the probed beliefs, not bounds.

\begin{proposition}[exact alpha-set values]\label{prop:pbvi}
The alpha-vector representation of the underlying theory (its
class-restricted rational alpha-vectors) evaluated at probed beliefs
returns the exact value: for horizons \(k = 1, 2, 3, 4\), three stock
levels, and seven random rational beliefs, the maximum over the
deduplicated alpha-set equals the maximum over definition-level
trajectory enumeration --- all \(84\) pairings equal. The four-step
alpha-set deduplicates \(17^{4} = 83{,}521\) action sequences to at
most \(545\) distinct outcome vectors per level (and to a single
vector at the top level's one-step horizon, where every action is
safe --- the silent region of the cube). Point-based selection
chooses where to evaluate, never what the value is.
\end{proposition}

\begin{proof}
Both sides are maxima of the same linear functional over the same
finite set of action sequences: the alpha-set is the image of the
sequence set under the map that records, for each parameter cell,
whether the sequence survives at the level; the maximum of a function
over a finite set equals the maximum over its image, and
deduplication removes only repeated vectors, never the attaining one.
Exactness of each entry is the definition-level simulation; no
discretization enters.
\end{proof}

\subsection{The antichain census}\label{census}

\begin{proposition}[the antichain, not the lattice]\label{prop:census}
The raw subset lattice carries \(2^{16} = 65{,}536\) subsets per
level --- \(16 \times 65{,}536 = 1{,}048{,}576\) over the grid ---
while the stored object is the antichain of maximal survivable sets --- the zero-suppressed
discipline of representing combinatorial set families by their maximal
elements (Minato, 1993):
\(16\) sets at the edge and \(32\) per level above it, \(496\) stored
sets in total. The per-level compression ratio is \(4{,}096\times\) at
the edge and \(2{,}048\times\) above it; and by
Lemma~\ref{lem:triangle} nothing is lost --- the antichain is the
whole classification, not a summarization of one.
\end{proposition}

\begin{proof}
The raw count is \(16 \times (2^{16} - 1) = 1{,}048{,}560\) nonempty
subsets; the headline uses the full lattice
\(16 \times 2^{16} = 1{,}048{,}576\), the larger and therefore
conservative figure. The stored count is
\(16 + 15 \times 32 = 496\) by Proposition~\ref{prop:bands}. The
antichain is lossless by the theory's kernel formula: the
value of any prior is the maximal survivable-set mass, read off the
stored sets alone.
\end{proof}

\subsection{The audited structure at a glance}\label{glance}

\begin{table}[t]
\centering\footnotesize
\caption{The cube's audited structure, one row per layer; every entry is a
verification-script output in exact rational arithmetic.}
\label{tab:glance}
\begin{tabular}{@{}l p{4.9cm}@{}}
\toprule
Layer & Exact statement \\
\midrule
Per-step drift & matched $+\tfrac3{10}$; $k$ mismatches: $-\tfrac12-\tfrac{2k}{5}$; hold $-\tfrac12$ \\
Cube accounting & drifts sum to $-8$ per step over the 16 cells \\
Singletons & $16$ survivable (exactly the floor's edge) \\
Pairs & exactly the $32$ Hamming-adjacent pairs; no other pair \\
Larger sets & none (three-set and four-set candidates fail) \\
Raw vs stored & $2^{16} = 65{,}536$ raw subsets per level, $496$ stored maximal sets \\
Observation ladder & ten probe schedules, values doubled at each rung \\
Deadline law & $z_0 \ge 1 + T/2$, $T = 0, \dots, 4$ (all cells) \\
\bottomrule
\end{tabular}
\end{table}

\subsection{A deadline instance}\label{deadline}

\begin{proposition}[the deadline law on the cube]\label{prop:deadline}
Hold for \(T\) steps (drift \(-\tfrac12\)), let the parameter be
revealed exactly at \(T\), then play matched forever (drift
\(+\tfrac3{10}\) on every branch). The companion theory's additive
deadline law applies with blind-drift cost \(d_{0} = \tfrac12\) and
zero probe excursion: viability holds exactly when
\[
z_{0} \;\ge\; 1 \;+\; \tfrac{T}{2}, \qquad T = 0, \dots, 4,
\]
verified by direct simulation on the grid and, independently, by a
mechanized proof (Section~\ref{methods}). The excursion vanishes
because the learned act rises on every branch immediately --- the
law's \(e_{p}\) term prices the probe step the free revelation does
not need.
\end{proposition}

\begin{proof}
Necessity: under the hold every cell drifts \(-\tfrac12\), so at
revelation every branch sits at \(z_{0} - T/2\), and survival to
revelation forces \(z_{0} - T/2 \ge 1\). Sufficiency: the hold keeps
every branch at \(z_{0} - t/2 \ge 1\) for \(t \le T\) exactly
under the same inequality, and after revelation the matched action
rises at \(+\tfrac3{10}\) on every branch, so no later constraint
binds. This is the theory's additive deadline law at
\(d_{0} = \tfrac12\), \(e_{p} = 0\); the direct branchwise argument
needs no scalar-additivity hypothesis, and the simulation certifies
the identity cell by cell on the grid. The mechanized layer proves
the same threshold as an \emph{equivalence} --- viability holds iff
\(z_{0} \ge 1 + T/2\) --- over an arbitrary ordered field and with
\(T\) symbolic, so neither the grid nor the restriction
\(T \le 4\) is load-bearing for the argument.
\end{proof}

\medskip
\noindent\textbf{Reciprocal check (computational-certification
companion).} The deadline law's radius structure is certified
externally by the certification companion (Abaee, 2026, Computational viability certification): on this
instance's audited domain, a radius-\(\delta\) belief cell inherits
the verdict exactly for \(\delta \le z_{0} - (1 + T/2)\), with
\(\delta = 0\) recovering the stored per-state bound and the
boundary cells zero-slack; the joint check is archived with the
companion's verification record.

\subsection{Related work, and what this paper is}
\label{priorart}

This paper is a computational companion, and the combinatorics it leans on is classical.
Both facts are stated here rather than left for a reader to discover.

\subsubsection{The place of this paper in the series}

The title's ``Scale~II'' refers to a companion relation, not to a volume number. The
theory being computed here is the belief-state safety value \(V_k(b)\) --- the maximal
probability of remaining safe for \(k\) steps under the disturbance's worst case --- together
with its support identity and \(\alpha\)-vector structure; that theory, and the
Sperner-bounded antichain formula behind its witness census, is developed in the companion
paper (Abaee, 2026, Belief-state safety values for viability under incomplete observation;
in this set, Paper~2). The present paper takes that theory's objects into the
curse-of-dimensionality regime on one explicit instance and classifies the instance
completely in exact arithmetic. It proves no new result about the theory itself. A reader
wanting the theorems should read that paper; a reader wanting the instance classified, and
the cost of doing so exactly, should read this one.

\subsubsection{Antichains, Sperner theory and Dedekind's problem}

An antichain in the Boolean lattice is a family of pairwise incomparable subsets. Sperner's
theorem (Sperner, 1927) gives the maximum size of an antichain on an \(n\)-set as
\(\binom{n}{\lfloor n/2\rfloor}\); counting \emph{all} antichains is Dedekind's problem,
with the asymptotic count \(2^{\binom{n}{\lfloor n/2\rfloor}(1+o(1))}\) due to Kleitman
(1966) and exact values known only for small \(n\) (Stanley, 2013; Engel, 1997, for the
standard treatments). Hamming adjacency, and the coding-theoretic machinery built on it, is
likewise classical.

\textbf{This paper proves no new theorem about antichains.} The census of
Section~\ref{census} is a computation \emph{on one instance}: \(1{,}048{,}576\) raw subset
evaluations collapsing to \(496\) stored maximal sets. That figure is an output, and its
value is that it was obtained exactly and in full. The classical results are the scaffold,
not the contribution.

\subsubsection{Point-based and exact computation for partially observed models}

Point-based value iteration and its successors (Pineau, Gordon and Thrun, 2003; Shani,
Pineau and Kaplow, 2013, for the survey) approximate belief-space value functions by
propagating a finite set of belief points, with error bounds in the usual formulations. The
distinction here is not the idea of a point set but the arithmetic: every quantity in this
paper is an exact integer or rational, so the catalogue is not an approximation with a
bound but a classification of the instance. That costs feasibility at scale --- which is why
the instance is a four-parameter cube and not a larger one, and why
Section~\ref{scope} says so. On the qualitative side, the complexity and decidability
landscape for partially observed models under qualitative objectives is established
(Chatterjee, Doyen and Henzinger, 2009); the companion paper's position relative to that
literature is stated there, and this paper adds nothing to it.

\subsubsection{What is claimed}

Given the above, the contribution is exactly four things.

\begin{enumerate}
\item \textbf{A complete classification of the four-parameter cube instance} --- every
  maximal blind survivable set, every band, and the observation ladder, in exact integer
  and rational arithmetic, with no approximation and no sampling.
\item \textbf{The cost figures for doing so.} The census collapse
  (\(1{,}048{,}576 \to 496\)), the four-step \(\alpha\)-set deduplication
  (\(83{,}521 \to 545\)), and the doubling of the observation ladder. These are the
  empirical content of the paper: they say what exactness costs at this size.
\item \textbf{The pairwise-Hamming structural restrictions as they apply here} --- that
  every blind survivable set consists of Hamming-adjacent cells, and that no three cells of
  the four-cube are pairwise adjacent, so no blind policy of any period keeps three or more
  cells safe. The adjacency notion is classical; its consequence for this instance is the
  result.
\item \textbf{Machine-checking of the layer}, described next.
\end{enumerate}

\subsection{Mechanized verification}
\label{mech}

The results above are additionally machine-checked. Sixteen modules of the project's
Lean~4 formalization are devoted to this paper's objects, in namespace
\texttt{Formalizations.EBC}: \texttt{EBC\_Bands}, \texttt{EBC\_Classification},
\texttt{EBC\_Deadline}, \texttt{EBC\_Dynamics}, \texttt{EBC\_ExactBelief},
\texttt{EBC\_Hamming}, \texttt{EBC\_Ladder} and \texttt{EBC\_Pairs}, with successive
versions of several. Together they carry \(134\) theorem and lemma declarations. The
mapping onto the sections here is by name: \texttt{EBC\_Hamming} carries the
pairwise-Hamming classification of Section~\ref{hamming}, \texttt{EBC\_Bands} the exact
bands, \texttt{EBC\_Ladder} the doubling of Section~\ref{ladder},
\texttt{EBC\_ExactBelief} the point-based evaluation of Section~\ref{pbvi}, and
\texttt{EBC\_Deadline} the deadline instance of Section~\ref{deadline}.

\noindent\textbf{No admitted gaps.} The full project builds cleanly under Lean~4 (\(60\) of \(60\) build jobs). In Lean an unfinished proof surfaces as the axiom
\texttt{sorryAx} in the axiom footprint of every declaration depending on it; a
source-level scan with comments stripped finds no \texttt{sorry}, \texttt{admit} or
\texttt{axiom} declaration anywhere in the formalization layer, all sixteen EBC modules
included. Declarations using classical reasoning depend only on \texttt{Classical.choice}, \texttt{propext} and \texttt{Quot.sound}, the standard axioms of Lean's classical logic.
\noindent\textbf{Verification provenance.} The project is pinned to Lean~4 in
\texttt{lean-toolchain} (currently \texttt{v4.34.1}). The mechanized checks reported in this
section were run under \texttt{v4.14.0}. On 2026-09-30 the source-level invariants were
re-verified against the current tree: a comment-stripped scan of all \texttt{.lean} files in
the project finds no \texttt{sorry}, \texttt{admit}, \texttt{axiom} or \texttt{constant}
declaration anywhere, and the declaration counts and named declarations quoted below were
re-counted and confirmed present. \textbf{The build has not been re-run since the toolchain
pin moved}, so the build-job figures below date from the \texttt{v4.14.0} run and should be
re-confirmed before submission.


\noindent\textbf{Scope.} What is machine-checked is the discrete combinatorial core --- the
structure the enumeration walks. The enumeration itself, including the census of
Section~\ref{census}, is a computation performed by the deposited scripts and is
\emph{not} formalized.

\subsection{Verification methods}\label{methods}

Every number is regenerated by one deterministic script in exact integer and rational arithmetic, standard library only. The battery opens with the pair-sum bound over all \(120\) pairs and \(17\) actions; the
adjacency-triangle check over all \(560\) triples; the exhaustive
one-, two-, three-periodic classification with \(60\)-step
certificates; the ten ladder values by survivable-subset enumeration
within each support; the \(84\) point-based pairings against
definition-level enumeration; the census counts; the deadline
thresholds by direct simulation; the dimension-scope battery (the
pair-sum bound at \(m = 5\) over all \(496\) pairs and \(33\)
actions; the ball radii \(r^{*}(m)\) for \(m = 3, \dots, 8\) with
the drift sign pattern and \(60\)-step simulations of the radius-one
ball at \(m = 5\) and the \(29\)-cell radius-two ball at
\(m = 7\)); and every headline number asserted present in this
text.

\textbf{A second, independent apparatus.} The script above enumerates;
it does not prove. Alongside it, the results of
Sections~\ref{hamming}--\ref{deadline} are formalized in a mechanized
layer (Lean~4, 60 modules under \texttt{Formalizations/}) with no axioms, no admitted
gaps, and no reliance on a computer-algebra oracle: every statement is
derived from the ordered-field interface alone, so it holds in any
model of that interface. The two apparatus were built separately and
agree where they overlap --- the ladder's first two entries
(\(1/16\) at the edge, \(1/8\) above it), the
adjacency-triangle fact behind the \(560\)-triple check, and the
\(16\)-singleton / \(32\)-pair census. Agreement of an enumeration
with a proof is evidence about both; it is not a substitute for either.

What each result now rests on, stated plainly:

\begin{itemize}
\item \textbf{Proved.} The band structure at the floor and above the
edge (Proposition~\ref{prop:bands}); the ladder's value form
(Section~\ref{ladder}); the deadline threshold
(Section~\ref{deadline}). These hold as stated over an arbitrary
ordered field, with the support size and the horizon symbolic.

\item \textbf{Enumerated, not proved.} The census, the \(84\)
point-based pairings, the \(83{,}521 \to 545\) deduplication, and the
dimension-scope battery at \(m \ne 4\). These are finite checks over
the cube; their content is the count, and a count is what the script
supplies.

\item \textbf{Simulated.} The \(60\)-step ball certificates at
\(m = 5\) and \(m = 7\), and the grid certification of the deadline
thresholds. For the deadline the simulation is now redundant --- the
threshold is proved --- and is retained as an independent cross-check.
\end{itemize}

\noindent Where this text previously attributed a result to simulation
or enumeration alone, the attribution above governs.

\subsection{Delimitations}\label{scope}

The blind class is declared and load-bearing: the policy may not
depend on the belief, so a single action sequence governs every
parameter branch and the pair-sum argument applies; a
\(z\)-feedback policy could branch its actions across branches (the
branches' stock paths diverge) and is outside the classification, as in the class discipline of the companion theory
(Abaee, 2026, An obstruction calculus), where the blind and
observation-dependent recursions are shown to differ --- the
feedback family contains the blind one, strictly, under
hypotheses of the same shape as this instance's. That result is
proved \emph{in the companion's model}; the correspondence with
the cube here is one of form, not of type, so it motivates
restricting to the blind class rather than proving that
restriction necessary. It does, however, fix the meaning of the
open question recorded above: since the two disciplines differ
strictly somewhere, whether a non-blind discipline beats the
pairs \emph{on the cube} cannot be settled by transferring the
companion's strictness, and is left open. Drifts are
deterministic per branch; probes are exact and cost no stock beyond
the blind decline they interrupt. The pairs-only classification is scoped to
\(m = 4\): in dimension five the constant matched action already
keeps the six-cell radius-one ball alive from every level
(Remark~\ref{rem:scope}), and the maximal sets for \(m \ge 5\) are
not classified here. Whether some non-blind discipline
beats the pairs on the cube is not claimed either way. Scaling beyond
sixteen cells --- where even the antichain's input enumeration becomes
the bottleneck --- is the next step, for which the pairwise-Hamming
instrument and the ball construction are the argued generalization
targets.

\subsection{Conclusion}\label{conclusion}

The exact discipline survives the curse of dimensionality and the
cube repays it with structure a pencil would never find: Hamming
governance of survivable sets, a doubling observation ladder, a
silent region at the top of the stock grid, and a four-order
antichain census. The stored object is the antichain; the argued
instrument is the pairwise bound; and the two together are what make
\(1{,}048{,}576\) subset evaluations collapse to \(496\) sets without
loss.

\subsection*{References}
{\footnotesize%% ------------------------------------------------------------------
%% Entries added 2026-09-29 to support subsection \ref{priorart}.
%% Chatterjee/Doyen/Henzinger copied verbatim from the sibling paper's
%% bibliography (already cited there), so it is consistent across the set.
%% NEEDING PUBLISHER VERIFICATION BEFORE SUBMISSION (pagination or exact
%% venue not confirmed here):
%%   Sperner 1927; Kleitman 1969; Pineau et al. 2003 (pages);
%%   Shani et al. 2013 (pages).
%% ------------------------------------------------------------------
Chatterjee, K., Doyen, L., Henzinger, T.A., 2009. Qualitative analysis of
partially-observable Markov decision processes. In: Int. Symp. Math. Found. Comput. Sci.
(MFCS 2009). LNCS, vol. 5734, pp. 635--646.
Engel, K., 1997. \emph{Sperner Theory}. Cambridge University Press.
Kleitman, D.J., 1969. On Dedekind's problem: the number of monotone Boolean functions.
\emph{Proc. Amer. Math. Soc.} \textbf{21}, 677--682.
Pineau, J., Gordon, G., Thrun, S., 2003. Point-based value iteration: an anytime algorithm
for POMDPs. \emph{Proc. IJCAI 2003}.
Shani, G., Pineau, J., Kaplow, R., 2013. A survey of point-based POMDP solvers.
\emph{Autonomous Agents and Multi-Agent Systems} \textbf{27}(1), 1--51.
Sperner, E., 1927. Ein Satz \"uber Untermengen einer endlichen Menge. \emph{Math. Z.}
\textbf{27}, 544--548.
Stanley, R.P., 2013. \emph{Topics in Algebraic Combinatorics}. Chapter 4, The Sperner
property. \url{https://math.mit.edu/~rstan/algcomb/algcomb.pdf}.
}


{\footnotesize
Abaee, A. (2026).. An obstruction calculus for viability under incomplete observation. Submitted.

Abaee, A. (2026).. Probabilistic sufficiency for the obstruction calculus: belief-state safety values under partial observation. Submitted.

Abaee, A. (2026). Computational viability certification under incomplete
observation: a continuous-to-finite bridge with exact certificates, a
solver-assisted certified campaign, and a finite-side reference library.
Submitted.

Minato, S. (1993). Zero-suppressed BDDs for set manipulation in combinatorial problems. In: Proceedings of the 30th Design Automation Conference (DAC '93), 272--277.


}

Lovejoy, W.S., 1991. Computationally feasible bounds for partially
observed Markov decision processes. Operations Research 39, 162--175.

\subsection*{Declarations}

\subsection*{Funding}
None.

\subsection*{Competing interests}
None.

\subsection*{Data availability}
No external data were used.

\subsection*{Code availability}
The verification script
\url{paper2_exact_belief_computation_v13_verification.py} is publicly
available at \url{https://github.com/MIKEAA2020/general-sustainability}
(folder \texttt{arena agent 1/paper rewrites/latex}); it reproduces all bounds, bands, ladder values, pairings,
census counts, and thresholds verbatim.

\subsection*{AI declaration}
GLM (Z.ai), Qwen (Alibaba Cloud) and DeepSeek AI assisted with drafting
and iterative review. The author reviewed and edited outputs and takes
responsibility for the final work.

\end{document}
